B.17's 50-state parameter sweep established that redistricting algorithm
parameters are insensitive: across five key parameters, the maximum
variation in expected Democratic seats nationally is $\leq 0.3$.
This paper reframes that result adversarially.
The question is not ``what happens when parameters vary?'' but
``what is the maximum partisan shift a hostile actor can achieve by
optimizing parameters after observing partisan implications?''

We define the most Republican-favorable parameter configuration
($\ufactor = 5\%$, $\acounty = 10$, $T = 1{,}000$, $\aswing = 1.20$,
$\wvra = 0.6$) and the most Democratic-favorable configuration
($\ufactor = 0.1\%$, $\acounty = 1$, $T = 100$, $\aswing = 1.05$,
$\wvra = 0.2$), and report the partisan outcome difference between the
two extremes.

The main result: the most Republican-favorable parameters produce
$\Dnat = 210.7$ expected Democratic seats; the most Democratic-favorable
parameters produce $\Dnat = 211.0$ seats.
The difference is 0.3 seats nationally---well within the B.7 seed-variance
noise floor (one standard deviation $\sigma \approx 0.8$ seats).

The \dia{}'s parameter pre-registration protocol eliminates this gaming vector
entirely: parameters are fixed in a SHA-256 formula before Census P.L.\ 94-171
data are released, preventing post-data optimization.
Even if pre-registration were absent, the 0.3-seat maximum effect is smaller
than any individual state's seat-level uncertainty, making parameter-based
gerrymandering both non-manipulative and non-attributable.
